% knowledge_graph_figure_vertical.tex
% Include this file in your main document with: \input{knowledge_graph_figure_vertical.tex}
% Required packages in main document: tikz, amsmath, subcaption, array
\begin{figure}[h]
\centering

% Top part - Factual triples
\begin{subfigure}[t]{\textwidth}
\centering
\small
\begin{tabular}{@{} >{\ttfamily}l @{\hspace{0.5em}} l @{}}
( & Dr. Marie Curie, \textbf{BornIn}, Poland )\\
( & Dr. Marie Curie, \textbf{DaughterOf}, Władysław Skłodowski )\\
( & Dr. Marie Curie, \textbf{StudiedAt}, University of Paris )\\
( & Dr. Marie Curie, \textbf{ReceivedAward}, Nobel Prize in Physics )\\
( & Dr. Marie Curie, \textbf{ReceivedAward}, Nobel Prize in Chemistry )\\
( & Dr. Marie Curie, \textbf{ResearchedIn}, Radioactivity )\\
( & Nobel Prize in Physics, \textbf{CategoryOf}, Physics )\\
( & Radioactivity Theory, \textbf{FieldWithin}, Physics )\\
( & Dr. Marie Curie, \textbf{CollaboratedWith}, Pierre Curie )\\
( & Pierre Curie, \textbf{AffiliatedWith}, University of Paris )\\
( & Radioactivity Theory, \textbf{ContributedBy}, Dr. Marie Curie )\\
( & Dr. Marie Curie, \textbf{DaughterOf}, Irène Joliot-Curie )
\end{tabular}
\caption{Facts as triples consisting of a head and tail entity and a relation between them.}
\label{fig:knowledge_graph_triples_vert}
\end{subfigure}

\vspace{1em} % Add some vertical space between the subfigures

% Bottom part - Knowledge graph
\begin{subfigure}[t]{\textwidth}
\centering
\begin{tikzpicture}[
    node/.style={draw, rounded corners, align=center, font=\footnotesize},
    person/.style={node, fill=cyan!25, ellipse},
    location/.style={node, fill=lime!30, ellipse},
    field/.style={node, fill=orange!25, ellipse},
    recognition/.style={node, fill=violet!20, ellipse},
    organization/.style={node, fill=lime!30, ellipse},
    edge/.style={->, font=\tiny, above}
]
% Nodes
\node[person] (marie) at (0, 0) {Dr. Marie\\Curie};
\node[person] (irene) at (-2.8, 2.2) {Irène Joliot-\\Curie};
\node[person] (father) at (-2.8, -1.8) {Władysław\\Skłodowski};
\node[location] (poland) at (-3.2, 0.3) {Poland};
\node[organization] (sorbonne) at (-1.2, -3.8) {University\\of Paris};
\node[person] (pierre) at (2.2, -3.8) {Pierre\\Curie};
\node[field] (radioactivity) at (1.2, 2.2) {Radioactivity\\Theory};
\node[field] (physics) at (4.2, 0.3) {Physics};
\node[recognition] (nobel_phys) at (3.8, -1.5) {Nobel Prize\\in Physics};
\node[recognition] (nobel_chem) at (4.0, -2.9) {Nobel Prize\\in Chemistry};
% Edges
\draw[edge] (marie) -- node[left] {\tiny DaughterOf} (irene);
\draw[edge] (marie) -- node[below] {\tiny DaughterOf} (father);
\draw[edge] (marie) -- node[above] {\tiny BornIn} (poland);
\draw[edge] (marie) -- node[below] {\tiny StudiedAt} (sorbonne);
\draw[edge] (marie) -- node[above left] {\tiny CollaboratedWith} (pierre);
\draw[edge] (pierre) -- node[below] {\tiny AffiliatedWith} (sorbonne);
\draw[edge] (radioactivity) -- node[right] {\tiny ContributedBy} (marie);
\draw[edge] (radioactivity) -- node[above] {\tiny FieldWithin} (physics);
\draw[edge] (marie) -- node[above] {\tiny ResearchedIn} (radioactivity);
\draw[edge] (marie) -- node[left] {\tiny ReceivedAward} (nobel_phys);
\draw[edge] (marie) -- node[left] {\tiny ReceivedAward} (nobel_chem);
\draw[edge] (nobel_phys) -- node[below] {\tiny CategoryOf} (physics);
\end{tikzpicture}
\caption{Corresponding knowledge graph visualization.}
\label{fig:knowledge_graph_visual_vert}
\end{subfigure}

\caption{Illustration of knowledge representation using factual triples (top) and their corresponding graph structure (bottom). Node colors represent entity types: \textcolor{cyan}{cyan} represents people, \textcolor{orange}{orange} represents research fields, \textcolor{violet}{violet} represents awards, and \textcolor{lime}{lime} represents locations and institutions.}
\label{fig:knowledge_graph_vert}
\end{figure}